\section{Related Work}\label{sec:related}

Classical time-series methods such as ARIMA and Holt-Winters handle linear temporal patterns reasonably well, but they consistently lose to machine learning once exogenous weather variables enter the picture. Almusaylh et al.\ \parencite{almusaylh2018forecasting} showed that SVR and MARS outperform ARIMA on short-term Australian demand, and Schirmer et al.\ \parencite{schirmer2019evaluation} found tree-based models clearly superior to linear regression for residential load prediction. Within the tree-ensemble family, Random Forest \parencite{breiman2001randomforest} controls variance via bagging; XGBoost \parencite{chen2016xgboost} introduces regularised sequential boosting and has become a go-to model in energy forecasting \parencite{khan2020hybrid}; and CatBoost \parencite{prokhorenkova2018catboost} offers ordered boosting with native categorical support---a useful property for demand data that includes day-of-week, month, and holiday flags. Weather, holidays, and weekends are well-established demand drivers \parencite{khan2021influencing}.

On the tuning side, meta-heuristic optimisation can make a substantial difference---Li et al.\ \parencite{li2020datadrivenstrategy} reported a 42\% accuracy gain from PSO-tuned models on university building data. PPSO \parencite{ghasemi2019ppso} is a variant of PSO that uses a phase angle to modulate the cognitive and social velocity components through periodic functions, which removes the need to hand-tune inertia and acceleration coefficients. Zhang et al.\ \parencite{zhang2023catboostppso} applied CatBoost-PPSO to hourly consumption of a single residential customer and found it outperformed five alternative meta-heuristics.

However, there is a methodological concern in prior work: tuning is typically applied only to the authors' preferred model, so it is unclear whether the reported advantage comes from the algorithm itself or simply from having better hyperparameters. Additionally, CatBoost-PPSO has not been tested on daily national-level demand in a tropical climate. Our project tackles both of these gaps.
